%!TEX program = uplatex
\RequirePackage{plautopatch}
\documentclass[uplatex,dvipdfmx]{jlreq}
\usepackage{amsmath,amssymb,amsthm}

\pagestyle{empty}

\newcommand{\CC}{\widehat{\mathbb{C}}}

\begin{document}

%======================================================================
% 表紙
%======================================================================
\begin{center}
  \vspace*{20pt}
  {\Huge\bfseries 複素力学系}

  \vspace{1pc}
  {\Large 概要}
\end{center}

\vfill

\newpage

%======================================================================
% 講演１：宍倉光広「複素力学系」（3回）
%======================================================================
\noindent
{\large\bfseries 複素力学系}\\
\hspace{2em} --- ジュリア集合とは ---\\
\hspace{2em} --- サリバンの非遊走定理と擬等角写像 ---\\
\hspace{2em} --- ２次多項式とその周辺 ---

\bigskip
\hspace*{\fill}{\large\bfseries 宍倉　光広（京都大学大学院理学研究科）}

\bigskip

複素力学系とは、１次元または高次元の複素多様体の上で定義された複素解析的な力学系を意味する。
ここでは特に、Riemann 球面上の有理関数の反復合成（iteration）について考える。
興味の対象は、反復合成によって定義される軌道の様子や不変集合の構造、そして、
パラメータを変化させたときのこれらの変化の様子、
さらには有理関数全体の空間の力学系的性質による分類などである。
複素力学系の研究は1910年代に Fatou や Julia 等により始められたが、
やがて下火になり、その後 Siegel や Baker 等の散発的結果はあったが、
あまり注目を集める分野ではなかった。
しかし、1980年代初頭から再び注目を集めるようになってきた。
その理由の一つには、コンピュータの発展により「カオス」的な力学系や
「フラクタル」集合が目に見える対象として提示され、
それらの複雑さ、多様性、おもしろさ、さらには美しさが身近になってきたということがあげられる。
もう一方の理由としては、擬等角写像や Teichmüller 空間論などの重要な道具が整備され、
画期的な結果が数多く得られたということである。
この講演では、３回に分けて、複素力学系の諸性質の概観と、
いくつかの中心的課題について述べていきたい。

$\CC$ を Riemann 球面とし、$f : \CC \to \CC$ を次数が２以上の有理関数とする。
$f$ の $n$ 回合成を $f^n$ で表す。$f$ の Fatou 集合 $F_f$ を
\[
  F_f = \left\{ z \in \CC \;\middle|\;
  z \text{ のある近傍 } U \text{ で }
  \{f^n\}_{n=0}^\infty \text{ が同等連続（あるいは正規族）}
  \right\}
\]
で定義し、その補集合 $J_f = \CC \setminus F_f$ を Julia 集合と呼ぶ。
Fatou 集合は開集合、Julia 集合は閉集合で、ともに $f$ に関して不変な集合となる。
Sullivan の非遊走定理によると、Fatou 集合の各連結成分（Fatou 成分）は $f$ で
何回かうつせば、最終的には（集合として）周期的になる。
さらに、周期的な Fatou 成分上では、軌道は吸引的または放物型の周期点に収束するか、
円盤または円環上の無理数回転に共役になる。
非遊走定理の証明では、擬等角写像とそれを用いた Riemann 面の等角構造の変形が、
決定的な役割を演じた。
これ以降、１次元複素力学系の研究では、
擬等角写像やその周辺の考え方が様々な局面で現れてくることになる。
また、この視点から複素力学系の理論と Klein 群の理論の様々な側面での類似
（Sullivan の辞書）は、双方の理論の研究の上で重要な指導原理となった。
（松崎氏の講演を参照。）

複素力学系の最も簡単な例は２次多項式である。
コンピュータによる数値実験では、
２次多項式に限っただけでも、非常に多様な形と豊富な分岐現象が起こりうることがわかる。
２次多項式に関する精力的な研究を始めたのは Douady と Hubbard であった。
external ray、external angle を導入し、ポテンシャル論と組み合わせて、
多項式の Julia 集合を位相的、組み合わせ的に記述をすることが出来るようになった。
Mandelbrot 集合は、２次多項式のパラメータ空間内で定義され、分岐を記述する集合であるが、
Douady と Hubbard はこの集合の組み合わせ的記述に成功している。

２次多項式に限っても、双曲型系の稠密性問題、局所連結性の問題、くりこみの問題など、
興味深い問題が関わってくる。
例えば、実２次多項式族のエントロピーの単調性問題などは、
問題を複素化してはじめて解決された。（辻井氏の講演を参照。）
時間の許す限りこれらの問題についても解説したい。

\vspace{1pc}

\begin{center}
  \textsc{参考文献}
\end{center}
\begin{enumerate}
  \item[{[B]}] A.~Beardon,
        \textit{Iteration of Rational Functions: Complex Analytic Dynamical Systems},
        Graduate Texts in Mathematics \textbf{132},
        Springer-Verlag, New York, 1991.
  \item[{[M]}] J.~Milnor,
        \textit{Dynamics in One Complex Variable: Introductory Lectures},
        Friedr.~Vieweg \& Sohn, Braunschweig, 1999.
  \item[{[UTM]}] 上田・谷口・諸澤，
        複素力学系序説，
        培風館, 1995.
\end{enumerate}

\newpage

%======================================================================
% 講演２：松崎克彦「クライン群のタイヒミュラー空間」
%======================================================================
\noindent
{\large\bfseries クライン群のタイヒミュラー空間}

\bigskip
\hspace*{\fill}{\large\bfseries 松崎　克彦（お茶の水女子大・理）}

\bigskip

クライン群とは、リーマン球面の双正則自己同相写像（メビウス変換）からなる離散群のことである。
1900年前後、Klein は Fricke とともにクライン群の力学系の研究を始め、
軌道の集積点集合である極限集合に注目した。
極限集合は、その後の Julia や Fatou による有理写像の反復合成の力学系の理論では
Julia 集合（作用がカオス的である集合）に相当するものであり、
当時からこの２つの複素力学系の類似についての認識はあったと思われる。
有理写像に比べてクライン群が理解しやすい点は、
対応する幾何学的対象が自然に存在するところである。
メビウス変換は上半空間に拡張し、双曲計量に関する等距離変換として作用する。
従って、その商空間としてあらわれる完備３次元双曲多様体がクライン群作用の幾何学化である。
この視点は Poincaré に源を発し、現代では Thurston の３次元双曲多様体論となっている。

リーマン面のモジュライ問題がタイヒミュラー空間の理論として発展した過程において、
擬等角写像の果たした役割は極めて大きい。
リーマン面の変形空間として素朴に考えられる基本群のホロノミー表現空間に比べて、
より内在的なパラメーター空間がこれにより与えられることになった。
この方法はクライン群の変形理論にも適用され、
完備３次元双曲多様体の無限遠に位置するリーマン面、
それはクライン群の作用が安定的である領域（極限集合の補集合で不連続領域や
Fatou 集合とも呼ばれる）の幾何学化であるが、
そのタイヒミュラー空間がクライン群の不連続領域上での擬等角変形空間となる。
自然な流れとして、
基本群のホロノミー表現空間の次元とタイヒミュラー空間の次元の比較により、
有限生成クライン群の不連続領域上での作用の幾何学的な有限性が帰結される。
1960年代の Ahlfors の有限性定理は、
こうしてクライン群のタイヒミュラー空間論のひとつの成果を示したのであるが、
実はそれが、有理写像とクライン群の複素力学系の相互関連の研究の新たな出発点となったのである。
Sullivan は有理写像の力学系に擬等角変形の理論を導入し、
Ahlfors 有限性定理と同様な原理で Julia、Fatou の時代からの未解決問題であった
遊走領域の非存在に証明を与えた。
さらに有理写像の複素力学系の種々の概念、指数、定理がクライン群の場合の何に対応するかをみることにより、
より幾何学的視点からの有理写像の力学系の考察を提案している（Sullivan の辞書）。

この講演では、クライン群、擬等角写像、タイヒミュラー空間の基本的な説明ののち、
上記の Ahlfors 有限性定理からはじめて、
Sullivan の辞書のクライン群側の項目についての解説をする。
とくに有限性定理では扱われていない極限集合上での擬等角変形について、
Mostow 剛性と構造安定性との関連について取り上げる。
３次元閉双曲多様体の双曲構造が位相的同値類によってのみ決定されるというのが
Mostow 剛性のひとつの定式化ではあるが、その本質的部分は、
クライン群が極限集合上で擬等角変形を許容しないことを示すことにある。
可測リーマン写像定理によれば、これは群作用で不変な可測ベクトル場の非存在と同値である。
Sullivan はクライン群の作用がある種のエルゴード性をもつとき、
不変可測ベクトル場の非存在を証明し、
とくに Ahlfors 有限性定理と同様の議論から、
有限生成クライン群は極限集合上で擬等角変形をもたないことを示し、
Mostow 剛性定理の拡張を与えた。
さらにこの結果を用いて、
クライン群が摂動によりその代数構造が変わらないという構造安定性が、
力学系の双曲性と同値であることも証明した。
有理関数の力学系の理論との好対照がここにあらわれる。
ジュリア集合上で不変可測ベクトル場の非存在を証明することは、
有理関数の理論の中心となる未解決の問題である。
一方、構造安定性をもつ系が変形空間のなかで稠密に存在することは、
有理関数に対しては示されているが、クライン群論では未解決である。
講演の最後では、双曲多様体の ending lamination 予想とタイヒミュラー空間の境界に関する
Bers--Thurston 予想について述べ、この未解決問題の最近の進展についても触れる予定である。

\newpage

%======================================================================
% 講演３：辻井正人「２次関数の力学系」
%======================================================================
\noindent
{\large\bfseries ２次関数の力学系}

\bigskip
\hspace*{\fill}{\large\bfseries 辻井　正人（北海道大・理）}

\bigskip

実１次元力学系というのは、区間から区間への写像 $f$ の反復合成を力学系と考えたものです。
このような素朴な力学系の中にもカオスと呼ばれる複雑な軌道の挙動が現れることはよく知られています。
特に、ここでは２次関数 $Q_a(x) = a - x^2$ の力学系について考えます。
このような力学系は複素力学系の実軸への制限と見なすこともできて、
いわば実と複素の１次元力学系の交差点のようになっています。
ただ、それだけではありません。
２次関数の力学系の解析は過去20年間のカオス的な力学系の研究の殆ど中心的な部分を占めてきました。
ここで「殆ど」と書いたのは本当の中心ではないからです。
中心はホモクリニック接触とそれを含む分岐にあります。
ただ、これらの分岐についてはそれに近づくのでさえ容易ではないので
力学系の研究者の多くはそれを横目で見ながらその周りを何かのとっかかりを探して歩き回っている
というのが（私から見た）現状です。
そういった中で、２次関数の力学系は最も素直でしかも有望なとっかかりで、
多くの力学系の研究者が２次関数の力学系に力を注いだことは自然なことです。
私の講演はこの辺の事情から話を始めたいと思います。

２次関数の力学系については、研究が始まった当初から数値計算に基づいて次の３つの予想がありました。
\begin{itemize}
  \item $Q_a$ が安定周期点を持つようなパラメータの全体は開かつ稠密。
  \item $Q_a$ のエントロピーは $a$ について広義単調増加。（単調性予想）
  \item $Q_a$ が strange attractor を持つようなパラメータの全体はルベーグ測度が正。
\end{itemize}
これらの予想は現在では全て肯定的に解決していて定理になっています。
聞いた話によれば、これらの予想はこの順番に難しくなると思われていたということです。
しかし、面白いことに、これらの予想が解けた年代は順に90年代、80年代、70年代となっています。
また、最初と２番目の予想は複素力学系の理論によって解決されています。
講演では特に２番目の単調性予想について、
それがどのように複素力学系の理論を使って解かれたかを説明します。

\end{document}
